\section{Introduction}
\label{sec:introduction}

Supersymmetric matrix quantum mechanics appears in the regularization of
supermembranes~\cite[Section~3]{DeWitHoppeNicolai1988} and in the
matrix-theory conjecture of Banks, Fischler, Shenker, and
Susskind~\cite{BFSS1997}. The model studied here describes nine Hermitian
matrices and their fermionic partners. Its commuting configurations form
noncompact flat directions.
The D0-brane bound-state prediction of Witten~\cite[Section~4.2]{Witten1996BoundStates},
subsequently central to matrix theory, is that after removing the free
center of mass every finite matrix size has exactly one normalizable
state of zero energy. We prove this threshold-bound-state conjecture for the
\(\SU(N)\) model on its full configuration space.
In the matrix-theory interpretation, this relative bound state supplies
the internal state of a single supergraviton carrying \(N\) units of
longitudinal momentum~\cite[Section~5]{BFSS1997}.
The continuous spectrum starts at zero~\cite{DeWitLuscherNicolai1989},
so the proposed eigenstate lies at its threshold rather than below a
gapped continuum.

\subsection{The operator and the result}
\label{sec:model}

Fix an integer \(N\ge2\) and put \(d_N=N^2-1\).
Choose a traceless Hermitian basis \(T_A\), \(1\le A\le d_N\), with
\[
\Tr(T_AT_B)=\delta_{AB},\qquad
T_AT_B-T_BT_A=\ii f_{ABC}T_C.
\]
Repeated indices are summed unless stated otherwise.
The bosonic coordinates and momenta are
\[
X_i=x_i^AT_A,\qquad p_i^A=-\ii\frac{\partial}{\partial x_i^A},
\qquad 1\le i\le9.
\]
Let \(\mathcal F_N\) be an irreducible complex Clifford module with
self-adjoint generators \(\theta_\alpha^A\), \(1\le\alpha\le16\), satisfying
\[
\{\theta_\alpha^A,\theta_\beta^B\}
  =\delta_{\alpha\beta}\delta_{AB}.
\]
Choose real symmetric \(16\times16\) matrices \(\gamma^i\) with
\(\{\gamma^i,\gamma^j\}=2\delta_{ij}\Id\), and set
\(\gamma^{ij}=[\gamma^i,\gamma^j]/2\), where the bracket in this formula
is the ordinary matrix commutator. The sixteen supercharges on smooth
compactly supported functions are
\begin{equation}
Q_\alpha
 =p_i^A\gamma^i_{\alpha\beta}\theta_\beta^A
 +\frac12 f_{ABC}x_i^Bx_j^C
       \gamma^{ij}_{\alpha\beta}\theta_\beta^A.
\label{eq:model-charge}
\end{equation}
The physical Hilbert space is
\[
\mathscr H_N
 =\bigl[L^2(\R^{9d_N})\otimes\mathcal F_N\bigr]^{\SU(N)},
\]
with the simultaneous adjoint action on bosons and fermions. Let \(H_N\)
be the nonnegative self-adjoint operator associated with the closure of
the quadratic form
\begin{equation}
q_N(\Psi)=\frac1{16}\sum_{\alpha=1}^{16}\norm{Q_\alpha\Psi}^2
\label{eq:model-form}
\end{equation}
on the gauge-invariant compactly supported smooth core. This fixes the
coupling and the operator domain used throughout the paper.
We use the exterior grading for the polarization
$T=\gamma^1\gamma^2\gamma^3$, declaring its empty vector even and
choosing that vector to minimize $\ii\theta^tT\theta/2$.
Section~\ref{sec:algebra} describes the associated symmetry actions.

\begin{theorem}\label{thm:main}
For every integer \(N\ge2\),
\[
\dim\Ker H_N=1.
\]
The kernel consists of \(\Spin(9)\)-invariant states and is even in the
fermion grading just specified.
\end{theorem}

All square integrability in Theorem~\ref{thm:main} is on the whole
space \(\R^{9d_N}\). Constants in the proof may depend on the fixed
integer \(N\). The free \(\U(1)\) center-of-mass coordinates do not occur
in \(\mathscr H_N\).

\subsection{History and the analytic issue}

The supersymmetric index counts even and odd zero modes with opposite
signs~\cite{Witten1982Index}. In the present noncompact problem,
contributions from infinity are essential. Yi~\cite{Yi1997} computed the
$\SU(2)$ bulk term and proposed a free-particle argument for its boundary correction;
Sethi--Stern~\cite{SethiStern1998} obtained the two-particle index by
including its boundary term. Moore--Nekrasov--Shatashvili~\cite{MNS2000}
computed the matrix integrals entering the principal contribution.
Green--Gutperle~\cite[Section~5]{GreenGutperle1998} argued that corrections
from free identical clusters account for the divisor sum, assuming the
free-cluster boundary prescription extends to general $N$.
Konechny~\cite[Sections~3 and~5]{Konechny1999} derived a free leading
asymptotic Hamiltonian and used it to conclude that the full $\SU(N)$
index equals one. His expansion assumes that the eigenvalue differences
of one diagonalized matrix grow at the same order. Our proof will
control also the degenerating and nested cluster regions through
explicit norm estimates.

A signed count becomes a dimension count once odd normalizable ground
states are excluded. Sethi--Stern~\cite[Sections~3.2 and~4.2]{SethiStern2000} proved rotational
invariance of normalizable ground states of ten-dimensional reductions.
Hasler--Hoppe~\cite[Theorem~1(a) and Lemma~2(c)]{HaslerHoppe2002} gave a direct
rotation anticommutator and cutoff proof on the full physical Hilbert
space. In the fermion representation used here, singletness implies
even parity. Section~\ref{sec:algebra} reproduces these arguments in
our normalization and fixes the closed operator domains. The recent
lectures of Lin~\cite[Section~3.1]{Lin2026v3} describe the finite-$N$ spectral picture
and distinguish the two-particle uniqueness consequence from the general
conjecture.

Asymptotic wavefunction analysis supplies a second important line of
work. Halpern--Schwartz~\cite{HalpernSchwartz1998} developed a generalized
Born--Oppenheimer search for $\SU(2)$ candidates. Graf--Hoppe~\cite{GrafHoppe1998}
computed the leading and subleading $\SU(2)$ asymptotics, and
Fr\"ohlich--Graf--Hasler--Hoppe--Yau~\cite[Section~3]{FroehlichEtAl2000} classified
the corresponding invariant formal supercharge expansions at infinity. For general rank, Hasler--Hoppe~\cite[pp.~2--5]{HaslerHoppeFactorization2002}
conjugated the charge by the slice density, projected onto its transverse
oscillator ground state, and obtained the leading free Cartan charge.
Their cancellation of the gauge-spin and density terms is a direct
predecessor of the computation in Section~\ref{fast:section}.
Lin--Yin~\cite[Sections~3.2 and~4]{LinYin2015} developed a recursive cluster-factorization
proposal and its first inverse-oscillator correction. These asymptotic
constructions motivate the local objects used below; a global
normalizable state additionally requires control of every cluster scale
and of all limiting norm.

Mass deformation offers a finite count from which to approach that
global problem. Porrati--Rozenberg~\cite{PorratiRozenberg1998} formulated
a deformation criterion with explicit normalizability requirements on
weighted representatives. Kac--Smilga~\cite{KacSmilga2000} counted
mass-deformed vacua and identified persistence of normalizability at
zero mass as the remaining hypothesis. We use the completely confining
plane-wave deformation of Berenstein--Maldacena--Nastase~\cite{BMN2002}.
Its classical zero orbits are indexed by integer partitions of $N$.
Dasgupta--Sheikh-Jabbari--Van Raamsdonk~\cite{DasguptaSheikhJabbariVanRaamsdonk2002}
computed the quadratic fluctuation spectra around the reducible as well as
irreducible vacua. Section~\ref{sec:massive-index} rederives these
frequencies and proves the stabilizer, domain, and norm-exhaustion
statements needed for the index of our specified invariant charge.
Chang's refined BMN index~\cite[Section~2.2]{Chang2025} also records signed
contributions from states with nonzero angular momentum. The count below
is the Fredholm index of a real charge on a specified invariant subspace.

The operator domain must remain fixed throughout this comparison.
Boulton--Garc\'ia del Moral--Restuccia~\cite[Sections~5--6]{Boulton2021}
establish existence and uniqueness for a valley boundary-value problem
with admissible prescribed trace. The realization with zero boundary trace has trivial
kernel, while the solutions with prescribed nonzero trace need not be
killed by the supercharges. That problem
therefore concerns a different realization from
\eqref{eq:model-form}.

The task addressed here is the passage from a massive Fredholm count
to the complete $L^2$ kernel of the undeformed full-space operator.
The local oscillator projections and fast-inverse corrections are
proved with graph-norm bounds near arbitrary proper clusters.
An exterior estimate, established simultaneously with the kernel
statement, excludes norm at intermediate internal scales. Together
with massive confinement, it accounts for every limiting norm and
inner product. This is the information that permits the final
unsigned dimension comparison.

\subsection{The proof and its compactness mechanism}

Write $\mathcal K_n=\Ker H_n$ and set $\mathcal K_1=\C$. The proof compares
$\mathcal K_N$ with a massive problem whose index is the partition number
$p(N)$, the number of unordered integer partitions of $N$. Two different
uses of partitions enter this comparison. At small coupling, they label
the massive classical vacuum orbits. In the limiting cluster decomposition, a proper
partition $C=(n_1,\ldots,n_k)$, $k\ge2$, specifies the sizes of separated
commuting blocks. Their distinct centers $c_\alpha\in\R^9$ satisfy
$\sum_\alpha n_\alpha c_\alpha=0$. A \emph{cluster profile} is a
wavefunction of these relative centers, with their fermions and with
values in $\bigotimes_\alpha\mathcal K_{n_\alpha}$.

For the dimension comparison, we study profiles of normalized massive
states whose charge norms tend to zero; we call these null profiles.
The proper partitions will contribute $p(N)-1$ lines of null profiles, while
a full-collapse profile takes values in $\mathcal K_N$. The comparison
uses norm-preserving extraction for existence, and recovery of prescribed
profiles together with a massive spectral gap for uniqueness.

\paragraph{The massive count.}
We write $D_v(h,m)$ for the deformed charge paired with a real unit
spinor $v$, where $h>0$ multiplies the commutator term and $m\ge0$
multiplies the linear mass term. Thus
$D_v(1,0)=\sum_\alpha v_\alpha Q_\alpha$.
Section~\ref{sec:algebra} chooses a spinor $u$ and a connected fixing
torus $G$ so that, on gauge and $G$ invariants,
\[
 D_u(h,m)^2\ge\frac23\mathcal H(h,m),
\]
where $\mathcal H$ is the averaged squared charge. At fixed positive mass,
the restricted charge is Fredholm. Section~\ref{sec:massive-index}
computes its index as $p(N)$ in the limit $h\downarrow0$ and proves
constancy for all $h>0$.

We then study $h\to\infty$ at mass one. The unitary change of variables
$y=h^{1/3}x$ sends $D_v(h,1)$ to
$h^{1/3}D_v(1,h^{-2/3})$, so this limit is a massless limit at the original
coupling. Internal block states shrink on the physical scale $h^{-1/3}$,
while their relative centers remain on scale one.

\paragraph{Local profiles and their equations.}
Near distinct block centers, a spectral slice separates the block
variables $b$ from the off-block variables. After rescaling the latter,
Section~\ref{fast:section} gives the exact decomposition
\[
 D_v(h,m)=\sqrt h\,E_v(b)+D_v^s(h,m)+R_v.
\]
Here $E_v$ acts on the transverse variables and their fermions,
$D_v^s$ is the charge of the block variables, and $R_v$ collects the
remaining slice terms. On the fast space fixed
by the block-center gauge torus, $E_v$ has an even one-dimensional
kernel and a local oscillator gap. A first-order estimate then forces
bounded graph-norm sequences onto this line. In the scalar-block limit,
compressing the charge onto this line gives the ordinary slow charge:
the gauge-spin term cancels the derivative of
the slice density.

To pass the charge equation to a profile, we need more than concentration
of the state. Section~\ref{lift:section} constructs test states whose
charge outputs converge in norm. A correction obtained from the reduced
fast inverse cancels the perpendicular charge output of the leading
oscillator product. Duality then identifies the center equation of each
profile. Fixed physical cutoffs and $L^2$ concentration make these
constructions available without moment bounds on the internal kernels.

\paragraph{Capturing every concentration scale.}
Local profiles leave a possible loss of norm between the internal scale
$h^{-1/3}$ and a vanishing physical radius. Section~\ref{sec:profiles}
excludes that loss by proving, simultaneously with Theorem~\ref{thm:main},
the exterior estimate
\begin{equation}
 \norm{rD_u(h,0)F}\ge c_N\norm F,
 \qquad
 \supp F\subset\{h^{1/3}r>R_N\},\qquad r=|x|.
 \label{eq:intro-exterior}
\end{equation}
It applies to compactly supported gauge-invariant vectors in the charge
graph domain, without a rotation restriction. The
smaller-block estimates make proper-cluster profiles account for all
norm on a compact annulus. Their center spaces have dimensions
$9(k-1)\ge9$, so the Euclidean Hardy inequality gives the exterior
estimate at the current size. Only after this step do we extract the
whole-block profile at the origin. The induction therefore uses no
prior existence or dimension statement for $\mathcal K_N$.

Section~\ref{sec:counting} proves a separate massive annular estimate
that prevents escape to infinity. The resulting decomposition preserves
all limiting norms and inner products. Each proper null profile has the
unique even massive center-oscillator vacuum as its center factor. Induction
on the internal blocks and the permutation symmetries then give one
line for each of the $p(N)-1$ proper partitions. Every null profile is even:
Proposition~\ref{prop:singlet-parity} proves this for $\mathcal K_N$
before its dimension is known. Consequently odd massive states have a
uniform charge gap at large $h$. Spectral pairing gives the same gap
on the even complement of the massive kernel, and the index becomes
the exact massive dimension $p(N)$. Extracting these $p(N)$ directions
forces a nonzero whole-block profile. Conversely, recovering any finite
orthonormal whole-block list along with the $p(N)-1$ proper lines and projecting
through the gap bounds that list by one. This closes the induction.

Figure~\ref{fig:induction} records this order of dependencies.
\begin{figure}[!htbp]
\centering
\begin{tikzpicture}[
  box/.style={draw=blue!55!black,rounded corners=2pt,align=center,
    text width=7.5cm,inner sep=7pt,font=\small},
  side/.style={box,text width=4.8cm,draw=black!60,
    inner sep=5pt,font=\footnotesize},
  flow/.style={-{Stealth[length=2mm]},thick,draw=black!70},
  node distance=6mm and 9mm]
\node[box] (input) {Induction input for $2\le n<N$:\\
  $\dim\mathcal K_n=1$ and exterior estimate $E_n$};
\node[box,below=of input] (proper) {Proper-cluster profiles capture all norm\\
  on compact annuli; corrected tests give center charges};
\node[box,below=of proper,fill=blue!7] (ext) {Hardy inequality in dimension $9(k-1)\ge9$:\\
  establish the current exterior estimate $E_N$};
\node[box,below=of ext] (collapse) {Capture the full-collapse profile in $\mathcal K_N$\\
  kernel nontriviality and dimension remain unknown};
\node[box,below=of collapse] (exhaust) {Exhaustion of norms and inner products;\\
  every limiting null profile is even};
\node[box,below=of exhaust] (gap) {Odd gap and spectral pairing;\\
  massive kernel has dimension $p(N)$};
\node[box,below=of gap] (count) {$p(N)-1$ proper profile lines;\\
  extraction gives existence;\\recovery and the gap give uniqueness};
\node[box,below=of count,fill=blue!7] (output) {Induction output: $E_N$ and $\dim\mathcal K_N=1$};
\node[side,right=of proper] (local) {Local fast analysis and\\corrected test lifts\\
  (Sections~\ref{fast:section}--\ref{lift:section})};
\node[side,right=of collapse] (tight) {Massive annular confinement\\
  prevents escape to infinity\\(Proposition~\ref{prop:massive-tightness})};
\node[side,right=of exhaust] (parity) {Current-size singlet and parity\\
  (Proposition~\ref{prop:singlet-parity}),\\$\dim\mathcal K_N$ unrestricted};
\node[side,right=of gap] (index) {Independent massive\\Fredholm index $p(N)$\\
  (Proposition~\ref{prop:massive-index})};
\foreach \a/\b in {input/proper,proper/ext,ext/collapse,collapse/exhaust,
  exhaust/gap,gap/count,count/output}\draw[flow] (\a)--(\b);
\draw[flow] (local)--(proper);
\draw[flow] (local)--(tight);
\draw[flow] (tight.west)--++(-0.3,0)|-(exhaust.east);
\draw[flow] (parity)--(exhaust);
\draw[flow] (index)--(gap);
\end{tikzpicture}
\caption{The simultaneous induction. Here $\mathcal K_n=\ker H_n$ and
$E_n$ is Proposition~\ref{prop:exterior-estimate} at size $n$.
The highlighted exterior step is completed before full collapse is
used. Massive confinement follows independently from the local
growing-mass estimates, and the index is computed independently at
small coupling. At $N=2$ all proper blocks are scalar, so the induction
input is empty.}
\label{fig:induction}
\end{figure}

\clearpage
